\begin{table*}[]
\vspace{4pt}
\resizebox{\textwidth}{!}{
\centering
\begin{tabular}{c|c|c|c|c|c|c|c|}
\cline{2-8}
 &
  \multicolumn{6}{c|}{\textbf{Task Time [seconds]}}&\multicolumn{1}{c|}{\textbf{Success Rate}}\\ \hline
\multicolumn{1}{|c|}{\textbf{Planner}} &
  \textbf{S1} &
  \textbf{S2} &
  \textbf{S3} &
  \textbf{S4} &
  \textbf{S5} &
  \textbf{S6} 
  & \textbf{S1 - S6}\\ \hline
\multicolumn{1}{|c|}{PokeRRT*} & 
\textbf{46.43 (21.67)} & 
212.74 (100.01) & 
232.42 (118.63) & 
70.47 (38.82) & 
132.62 (98.27) & 
\textbf{130.01 (84.05)} & 
0.87 (0.31) \\ \hline
\multicolumn{1}{|c|}{PokeRRT} & 
49.69 (21.11) & 
\textbf{196.54 (124.61)} & 
\textbf{167.55 (138.63)} & 
\textbf{64.14 (52.45)} & 
\textbf{116.35 (89.72)}& 
171.77 (103.21) & 
\textbf{0.88 (0.31)} \\ \hline
\multicolumn{1}{|c|}{\begin{tabular}[c]{@{}c@{}}LI-PokeRRT*\end{tabular}} & 
169.87 (71.28) & 
378.90 (122.21) & 
346.70 (111.71) & 
165.91 (44.87) & 
N/A & 
N/A & 
0.53 (0.28) \\ \hline
\multicolumn{1}{|c|}{\begin{tabular}[c]{@{}c@{}}LI-PokeRRT\end{tabular}} & 
123.71 (36.72) & 
328.66 (107.18) & 
322.28 (138.05) &
135.80 (32.91) & 
N/A & 
N/A & 
0.63 (0.18) \\ \hline
\multicolumn{1}{|c|}{\begin{tabular}[c]{@{}c@{}}Push Planner\end{tabular}} & 
122.68 (47.07) & 
284.78 (104.40) & 
249.32 (68.30) & 
117.58 (67.05) & 
N/A & 
N/A & 
0.44 (0.26) \\ \hline\hline
\multicolumn{1}{|c|}{Pick-and-Place} & 
16.29 (3.76) & 
17.71 (3.81) & 
16.14 (3.25) & 
N/A & 
19.15 (4.99) & 
N/A & 
0.67 (0.00)\\ \hline
\end{tabular}
}
\caption{Task times [\textsl{mean (stddev)}] for various planning algorithms in simulation are presented. Success rates are aggregated across all 6 scenarios. Overall, poke planning is faster than push planning and leads to higher success rate than pushing or grasping. Low-impulse (LI) poke planners and the push planner are unsuccessful in S5 and S6 due to action path collisions (S5) or goal region being outside the robot's reachable workspace (S6). The robot is unable to pick or place object in S4 and S6  due to limitations in robot mechanics.\vspace{-12pt}}
\label{tab:eval-sim-time}
\end{table*}